\SourcePage{209}{212}
\section{Boundary conditions in the semiclassical case}

Let $x=a$ be a turning point, so that $U(a)=E$, and suppose $U>E$ for every
$x>a$; the region to the right of the turning point is then classically
inaccessible. The wave function must decay into the interior of this region.
At a sufficient distance from the turning point it has the form
\begin{equation}
  \psi=\frac{C}{2\sqrt{|p|}}
  \exp\left(-\frac1\hbar\int_a^x|p|\,dx\right),\qquad x>a,
  \tag{47.1}
\end{equation}
corresponding to the first term of (46.10). To the left of the turning point,
the wave function must instead be represented by a real combination of the
two semiclassical solutions (46.9) of the Schrödinger equation:
\begin{equation}
  \psi=\frac{C_1}{\sqrt p}\exp\left(\frac i\hbar\int_a^x p\,dx\right)
  +\frac{C_2}{\sqrt p}\exp\left(-\frac i\hbar\int_a^x p\,dx\right),
  \qquad x<a.
  \tag{47.2}
\end{equation}

To determine the coefficients in this combination, one must trace the change
in the wave function from positive to negative values of $x-a$, beginning in
the region where (47.1) applies. This path, however, passes through the
neighborhood of the turning point, where the semiclassical approximation
fails and the exact solution of the Schrödinger equation must be considered.
For small $|x-a|$,
\begin{equation}
  E-U(x)\approx F_0(x-a),\qquad
  F_0=-\left.\frac{dU}{dx}\right|_{x=a}<0;
  \tag{47.3}
\end{equation}
in other words, this region presents the problem of motion in a constant
field. The exact solution of the Schrödinger equation for that problem was
found in §~24. The relation between the coefficients in (47.1) and (47.2) can
be found by comparing with the asymptotic expressions (24.5) and (24.6) of
that exact solution on the two sides of the turning point. It should be noted
that (47.3) gives $p(x)=\sqrt{2mF_0(x-a)}$, and therefore
\SourcePage{210}{213}
\[
  \frac1\hbar\int_a^x p\,dx
  =\frac{2}{3\hbar}\sqrt{2mF_0}(x-a)^{3/2};
\]
this coincides with the expression in the argument of the exponential or sine
in (24.5) or (24.6). The essential point in this reasoning is that the ranges
of validity of expansion (47.3) and of the semiclassical approximation partly
overlap. If the motion is semiclassical in almost the entire field region, as
has been assumed, there are values of $|x-a|$ small enough for (47.3) to be
valid and at the same time large enough to satisfy the semiclassical condition
and permit the asymptotic forms (24.5), (24.6).\footnote{Indeed, expansion
(47.3) applies for $|x-a|\ll L$, where $L$ is the characteristic distance over
which the field $U(x)$ changes. The semiclassical condition (46.7), on the
other hand, requires
$|x-a|^{3/2}\gg\hbar/\sqrt{m|F_0|}$. These requirements are compatible because
semiclassical motion far from the turning point (that is, for $|x-a|\sim L$)
means that $L^{3/2}\gg\hbar/\sqrt{m|F_0|}$.}

There is, however, another and methodically more instructive procedure that
avoids any need to invoke the exact solution. Formally regard $\psi(x)$ as a
function of the complex variable $x$, and pass from positive to negative
$x-a$ along a path lying wholly far from $x=a$, so that the semiclassical
condition is formally satisfied everywhere on the path (A.~Zwaan, 1929). We
again consider values of $|x-a|$ for which (47.3) is valid as well. The wave
function (47.1) then becomes
\begin{equation}
  \psi(x)=\frac{C}{2(2m|F_0|)^{1/4}(x-a)^{1/4}}
  \exp\left\{-\frac1\hbar\sqrt{2m|F_0|}
  \int_a^x\sqrt{x-a}\,dx\right\}.
  \tag{47.4}
\end{equation}

First trace how this function changes when the point $x=a$ is passed from
right to left along a semicircle of radius $\rho$ in the upper half of the
complex $x$ plane. On this semicircle
\[
  x-a=\rho e^{i\varphi},\qquad
  \int_a^x\sqrt{x-a}\,dx=\frac23\rho^{3/2}
  \left(\cos\frac{3\varphi}{2}+i\sin\frac{3\varphi}{2}\right),
\]
where the phase $\varphi$ changes from 0 to $\pi$. The modulus of the
exponential factor in (47.4) first increases (for
$0<\varphi<2\pi/3$), and then decreases to unity. At the end of the passage,
the exponent becomes purely imaginary and is equal to
\[
  -\frac i\hbar\int_a^x\sqrt{2m|F_0|(a-x)}\,dx
  =-\frac i\hbar\int_a^x p(x)\,dx.
\]
\SourcePage{211}{214}
Meanwhile, the prefactor in (47.4) changes under the continuation according to
\[
  (x-a)^{-1/4}\longrightarrow(a-x)^{-1/4}e^{-i\pi/4}.
\]
Thus the whole function (47.4) becomes the second term of (47.2), with
$C_2=(1/2)Ce^{-i\pi/4}$.

There is a simple explanation for why continuation through the upper half
plane determines only $C_2$ in (47.2). If the change of (47.2) is followed
along the same semicircle in the reverse direction, from left to right, its
first term rapidly becomes exponentially small compared with the second at
the beginning of the path. The semiclassical approximation cannot detect
exponentially small terms in $\psi$ against a large principal term; this is
why the first term in (47.2) is ``lost'' during this continuation.

To determine $C_1$, the point must instead be passed from right to left along
a semicircle in the lower half of the complex $x$ plane. The same reasoning
shows that (47.4) then becomes the first term in (47.2), with
$C_1=(1/2)Ce^{i\pi/4}$.

Consequently, the wave function (47.1) for $x>a$ corresponds for $x<a$ to
\[
  \psi=\frac C{\sqrt p}\cos\left(\frac1\hbar\int_a^x p\,dx+\frac\pi4\right).
\]
The resulting connection rule may be written in a form independent of which
side of the turning point contains the classically inaccessible region:
\begin{equation}
  \frac{2C}{\sqrt{|p|}}\exp\left\{-\frac1\hbar\int_a^x|p|\,dx\right\}
  \longrightarrow
  \frac C{\sqrt p}\cos\left\{\frac1\hbar\int_a^x p\,dx-\frac\pi4\right\}
  \quad
  \substack{U(x)>E\qquad U(x)<E}
  \tag{47.5}
\end{equation}
(H.~A. Kramers, 1926).
\SourcePage{212}{215}
We emphasize once more a point already evident from the derivation: this rule
is tied to a definite boundary condition imposed on one side of the turning
point, and in this sense it is to be used only in one direction. Specifically,
(47.5) was derived with the boundary condition $\psi\to0$ into the interior of
the classically inaccessible region; it must be applied in passing from that
region to the classically allowed region, as indicated by the arrow in
(47.5).\footnote{A passage in the reverse direction loses its meaning in the
sense that even a small change of the wave function on the right-hand side of
(47.5) may produce an exponentially growing term in the function on the
left-hand side.}

If the classically accessible region is bounded at $x=a$ by an infinitely
high potential wall, the boundary condition on the wave function at $x=a$ is
$\psi=0$ (see §~18). The semiclassical approximation then remains valid right
up to the wall, and the wave function is
\begin{equation}
  \begin{aligned}
    \psi&=\frac C{\sqrt p}\sin\left(\frac1\hbar\int_a^x p\,dx\right),
      &&x<a,\\
    \psi&=0,&&x>a.
  \end{aligned}
  \tag{47.6}
\end{equation}
